\section{Coarse-Graining II: Cross-Scale Shadows as Lagged Predictive Priors}\label{sec:cross_scale_shadows}

Coarse-graining I built a single meta-agent from a coherent cluster, the upward leg of the construction. This section supplies the downward leg: how the meta-agent, once formed, sets the priors of its constituents. In the construction the prior $p_i^{(\ell)}$ and the hyper-prior $r_i^{(\ell)}$ at level $\ell$ are not independent dynamical fields. They are the transported belief and the transported model of the meta-agent $I$ at level $\ell+1$ that contains agent $i$ as a constituent,
\begin{equation}
p_i^{(\ell)}(c) = \Omega_{iI}\big[q_I^{(\ell+1)}\big](c),
\qquad
r_i^{(\ell)}(c) = \tilde\Omega_{iI}\big[s_I^{(\ell+1)}\big](c),
\label{eq:cross_scale_shadow}
\end{equation}
where $\Omega_{iI}$ is the gauge-covariant transport into agent $i$'s frame on the latent-state fiber and $\tilde\Omega_{iI}$ is its model-fiber counterpart. The state transport is $\Omega_{iI} = U_i U_I^{-1}$ with $U = \exp(\phi)$ (Section~\ref{sec:transport}); the model transport is the $\rho_{\mathrm{model}}$-image of $U_i U_I^{-1}$ when frames are shared, or an independent transport $\tilde\Omega_{iI} = \tilde U_i \tilde U_I^{-1}$ with $\tilde U = \exp(\tilde\phi)$ when model frames are independent. Under this assignment beliefs and priors share the latent-state manifold $\mathcal{B}_{\mathrm{state}}$, models and hyper-priors share $\mathcal{B}_{\mathrm{model}}$, and the four-field hierarchy $r \to s \to p \to q$ posited at the agent definition reduces to two primitive fields $(q, s)$ at each level connected by cross-scale transport: $r$ regularizes $s$ because $r$ is the meta-agent's model transported down, and $p$ regularizes $q$ because $p$ is the meta-agent's belief transported down.

Read as an assignment at one instant, Eq.~\eqref{eq:cross_scale_shadow} sets a prior equal to a transported posterior of the same instant. That is not a factor of one normalized joint: a generative kernel that takes a recognition law as its argument is a map from recognition laws to kernels, and the free energy built from it is not an evidence bound~\cite[Chs.~4 and 12]{DennisGaugeVFE2026}. Appendix~\ref{app:augmented_joint} makes the failure concrete. There a joint with Gibbs cross-scale factors is a well-defined density on the product space of all agent-level latents (Lemma~\ref{lem:aug_joint_welldefined}), and Lemma~\ref{lem:shadow_mf_optimum} computes the rigid-link variational predictive density of a child under that joint, which retains the transported parent covariance and is not a message of the joint (Remark~\ref{rem:shadow_marginal_vs_mf}). Evaluated against the full parent posterior, the shadow reuses each constituent's own evidence: in the rigid-link limit the child precision obtained from shadow times likelihood exceeds the exact child posterior precision by exactly the precision of the child's own observation, which has entered the parent posterior and then enters the child a second time. For the same reason the same-time shadow functional is not bounded below by the surprisal of the augmented joint.

Read across time, Eq.~\eqref{eq:cross_scale_shadow} is exact. Index episodes by $n = 0, 1, \ldots$, and let the tower of episode $n+1$ be generated from the parent latent $z_I$ of episode $n$ by a declared transition kernel, with a declared prior on the tower of episode $0$. This episode-indexed joint is normalized and acyclic, every kernel reads only structural data and latents of earlier episodes, and the prior of the tower at episode $n+1$ given the records $o_{0:n}$ is the filtered parent law pushed through the transition~\cite[Ch.~12, the delayed-feedback proposition]{DennisGaugeVFE2026}. Choose the transported-copy transition, which places child $i$'s latent at episode $n+1$ at $\Omega_{iI} z_I$. The predictive prior of child $i$ is then
\begin{equation*}
p_i^{(\ell), n+1} = (\Omega_{iI})_{\#} \mathbb{P}\big(z_I^{(n)} \in \cdot \mid o_{0:n}\big),
\end{equation*}
which is $\Omega_{iI}[q_I^{(\ell+1)}]$ whenever $q_I^{(\ell+1)}$ is the filtered parent law of episode $n$, and is the corresponding variational predictive density when $q_I^{(\ell+1)}$ approximates it (status D; the display is the pushforward of the filtered law under a deterministic kernel). A noisy transition $\mathcal{N}(\Omega_{iI} z_I, \Omega_{iI} \Sigma_\epsilon \Omega_{iI}^\top)$ adds $\Omega_{iI}\Sigma_\epsilon\Omega_{iI}^\top$ to the transported parent covariance. No evidence is counted twice, because child $i$'s record of episode $n+1$ has not entered the filter that produced the parent law of episode $n$. Eq.~\eqref{eq:cross_scale_shadow} is therefore the one-episode-lagged predictive prior of the episode-indexed joint, and this lagged construction of the companion theory supersedes the same-time shadow. The same-time form used in the energy of the following sections is an energy-level simplification of the lagged scheme (status S), not a limit of it.

Two further readings are consistent with an evidence bound within one episode. The standard hierarchical scheme~\cite{friston2008hierarchical, parr2022active} derives the level-$\ell$ prior from a declared conditional $p(z_\ell \mid z_{\ell+1})$ of the generative model; its variational term is the conditional coupling $\mathbb{E}_{q_I} D_{\mathrm{KL}}(q_i \| \chi(\cdot \mid k_I))$ of a downward kernel $\chi$~\cite[Ch.~12]{DennisGaugeVFE2026}, and the corresponding downward message is the cavity predictive $\int \chi(\cdot \mid k_I) q_I^{\setminus i}(dk_I)$, which uses the parent posterior with constituent $i$'s contribution removed. The second is a random-effects model: the constituents are independent draws $k_i \sim \mathcal{N}(\Omega_{iI} m_I, \Omega_{iI} S_I \Omega_{iI}^\top)$ with a point hyperparameter $\theta_I = (m_I, S_I)$, so that $q_I$ is a point estimate, $S_I$ is the population covariance, and minimizing the free energy over $(q, \theta_I)$ is type-II maximum likelihood, whose exact M-step is the strict barycenter with its dispersion term (Section~\ref{sec:meta_agent_variational}). Appendix~\ref{app:augmented_joint}'s single parent latent with noisy children is a different model, in which $\Sigma_I$ is posterior uncertainty and not population spread; its rigid-link lemma illustrates the predictive density and does not supply the shadow as a factor.

The single rung extends to a tower. Meta-agents at level $\ell+1$ can themselves cohere into meta-meta-agents at level $\ell+2$, so that individuals form groups, groups form communities, and communities form societies up to a maximum level $\ell_{\max}$, with information flowing bottom-up through consensus formation and top-down through predictive priors. This bidirectional coupling has the feedback structure of Wheeler's participatory loop (W4 in Table~\ref{tab:wheeler_correspondence}; status S). That structure alone does not distinguish the construction from hierarchical modeling with learned priors~\cite[Ch.~14]{DennisGaugeVFE2026}; formation of collectives by the agents' own inference would, and Section~\ref{sec:participatory} states the test of that reading. The hierarchy has two boundaries. At the top level there is no enveloping meta-agent, and the hyper-prior $r_i^{(\ell_{\max})}$ is a fixed declared prior, the root of the hierarchy (Section~\ref{sec:ouroboros}); at the bottom level $\ell = 0$ there are no constituents, so cross-scale propagation acts only downward from above. A prospective single-level comparison may treat $p_i, r_i$ as primitive boundary conditions, while a multi-level run must test the cross-scale relation of Eq.~\eqref{eq:cross_scale_shadow}.

\section{Coarse-Graining III: Exact Coarse-Graining and Closure}\label{sec:rg_status}

The coarse-graining of the preceding sections is an exact block-variable pushforward of a posited configuration ensemble followed by a truncation to the five-term family, a block transformation in Kadanoff's sense~\cite{kadanoff1966scaling} without a rescaling step. It is motivated by, but is not, the renormalization group of statistical physics~\cite{Wilson1971, Wilson1974, Cardy1996}. No rescaling, coupling-space flow, or fixed point is constructed, and on a finite network none can be: a step that reduces the number of agents admits no comparison identification of the interaction spaces before and after it, and with blocking ratio at least $b > 1$ the tower has at most $\log_b N$ levels~\cite[Ch.~11, the finite-network obstruction proposition]{DennisGaugeVFE2026}. The partitions are chosen by a state-dependent detector, so the map is not a homogeneous flow, and the attention weights $\beta_{ij}$ are state variables (the softmax of the current edge energies), not coupling constants; the couplings proper, the temperature $\tau$, the weight $\lambda_h$, and the declared precisions, are held fixed across levels by assumption, and the exact pushforward would in general change them; the one exception computed here is the parent's entropy weight, which shifts at one loop (below).

The status of the construction rests on three claims. First, the coarse-graining map admits an exact pushforward, and it is gauge covariant (Theorem~\ref{thm:rg_covariance}) for any Lie group with a left-invariant Riemannian metric when the frames of each cluster lie inside the convexity radius; compactness enters only through the normalizability of the Haar reference of the Gibbs ensemble, which fails for noncompact $\mathrm{GL}^+(K_q)$. For an exact channel the fine free energy equals the coarse free energy plus the discarded conditional information, with the same evidence at both levels~\cite[Ch.~9, the exact coarse free-energy theorem]{DennisGaugeVFE2026}, and the coarse action is the conditional-partition formula of~\cite[Ch.~11]{DennisGaugeVFE2026}. Second, closure of the level-$(\ell+1)$ free energy in the same five-term family is conditional and posited, not derived. Third, the consensus detector is used as a candidate filter in place of the variational formation criterion; it is not derived from that criterion (Section~\ref{sec:meta_agent_threshold}). One coupling does change under the step. At one loop in the mean sector, at coherence, the fluctuation term $\tfrac{T_{\mathrm{cfg}}}{2}\log\det' H_I^{\perp}$ of the internal modes equals, up to a $Y$-independent constant, $-\tfrac{T_{\mathrm{cfg}}}{2}(|I|-1)\log\det\Sigma_I$, so integrating out the internal modes shifts the coefficient of $\log\det\Sigma_I$, the entropy weight of the parent, in proportion to the cluster size (Eq.~\eqref{eq:rg_entropy_shift}); this induced renormalization vanishes at $T_{\mathrm{cfg}} = 0$, where the step reduces to constrained minimization.

Several precedents organize coarse-graining in these terms. B\'eny and Osborne~\cite{beny2015information} formulate coarse-graining as a family of channels and pull the Fisher metric back onto the space of models, so that the metric contraction measures the information destroyed per step; Machta et al.~\cite{Machta2013} trace effective theories to the compression of Fisher-information eigenvalues under coarse-graining; Berman et al.~\cite{BermanKlingerStapleton2023} read Bayesian updating as a flow on the Fisher manifold; and Mehta and Schwab~\cite{mehta2014exact} construct an exact mapping between Kadanoff's variational renormalization group~\cite{Kadanoff1975, KadanoffHoughtonYalabik1976} and restricted Boltzmann machines. Exact form preservation of a coarse-graining step in the strict sense is the classic solvable-model fact (for example one-dimensional Ising decimation), not a claim we attribute to any of these references. The closure statement here is a parametric form preservation: the level-$(\ell+1)$ functional has the same five-term shape as the level-$\ell$ functional, with coarse belief couplings $\beta^{(\ell+1)} = f_\beta(\{\beta_{ij}^{(\ell)}\})$ pooled by the raw-conductance rule of Eq.~\eqref{eq:rg_inter_cluster}, in which the pooled intra-cluster conductance acts as a parent self-loop, and model couplings $\gamma^{(\ell+1)} = f_\gamma(\{\gamma_{ij}^{(\ell)}\})$ left unspecified; the working forms, the conductance pooling, and the Schur elimination behind them are in Appendix~\ref{sec:meta_agent_rg}. Lemma~\ref{thm:rg_exact_closure} (mean-sector Gaussian elimination) concerns only the linear coarse-graining of the means at fixed covariances, frames, and couplings, where the exact pushforward of the mean-sector Gibbs density is Gaussian with the Schur-complement precision. A map that retains the parent covariance has bounded fibers and lies outside that statement; a level-dependent log-determinant generates non-quadratic terms at the next level, so the class is not closed under iteration, and exact closure generically requires hyperedges~\cite[Ch.~11, the hyperedge-generation proposition]{DennisGaugeVFE2026}. A merely local quadratic expansion retains the anharmonic remainder and is only a Laplace approximation.

The belief-mean sector shows which operation preserves the family. There a functional of the five-term family is a block-diagonal positive part plus a matrix-weighted Laplacian with positive definite edge weights, so each off-diagonal block satisfies $\operatorname{sign}\det\Lambda_{ij} = (-1)^{K_q}$ and the triple product around any triangle carries the same sign, a necessary condition involving neither frames nor transport. Integrating the within-cluster deviations out at fixed barycenter violates this condition on sparse fine graphs: on a chain of pairs it generates a coupling between non-adjacent clusters with the wrong sign, the standard way a decimation leaves a nearest-neighbor ferromagnetic ansatz. The condition is only necessary, and it is weak at even $K_q$, so it certifies failure and not success. The general statement is in the companion supplement: unrestricted matrix-weighted Kron reduction leaves the Laplacian-plus-self-term family, and on a fixed flat graph closure for every weight assignment holds exactly when every hidden component has at most two retained terminals~\cite[Ch.~G7]{DennisGaugeVFESuppG2026}; the scalar-weight theory of Kron reduction and graph coarsening is in~\cite{DorflerBullo2013Kron, Loukas2019, villegas2023laplacian}. Identifying constituents preserves the family exactly. The congruence $S^\top\Lambda S$ is the Galerkin coarse form of aggregation-based algebraic multigrid; the within-cluster Laplacian contributions cancel identically, the coarse coupling between two clusters is the sum of the fine weights on the edges cut between them, which supplies $\beta^{(\ell+1)}$ directly without the conductance ansatz, and successive aggregations compose~\cite[Ch.~G9]{DennisGaugeVFESuppG2026}. Tying has two distinct probabilistic readings. As a coarse model it is conditioning on zero within-cluster disagreement, with natural parameters $(S^\top \mathsf{h}, S^\top\Lambda S)$~\cite[Ch.~G7]{DennisGaugeVFESuppG2026}; this is precomposition, which defines a new model to which data processing and evidence preservation do not transfer~\cite[Ch.~9]{DennisGaugeVFE2026}. As a restriction of the recognition family on the means it keeps the fine joint and costs bound. The coarse couplings $\beta^{(\ell+1)}$ above are those of the first reading. Exact marginalization, a Markov pushforward, moves neither the evidence nor the bound. The rigorous literature on position-space coarse-graining maps shows that block-spin and decimation images of Gibbs measures can fail to be Gibbsian~\cite{VanEnterFernandezSokal1993}; its comparison between schemes does not settle the direction the framework's intuition assumes, so no analogy in either direction is asserted, and the lattice treatment of spatial coarse-graining is deferred to~\cite[Chs.~B7--B9]{DennisGaugeVFESuppB2026}.

The construction quantifies cross-scale information propagation through the transported relative entropy between constituent and meta-agent beliefs,
\begin{equation}
\mathcal{I}_{\ell \to \ell+1} = \sum_{I \in \mathrm{level}(\ell+1)} \sum_{i \in I} D_{\mathrm{KL}}\big(q_i^{(\ell)} \big\| \Omega_{iI}\big[q_I^{(\ell+1)}\big]\big),
\label{eq:cross_scale_information_flow}
\end{equation}
a divergence between two distributions on the same fiber after transport of the meta-agent belief into the constituent's frame, which vanishes when they agree after transport. We do not write this quantity as a mutual information, since a mutual information requires a joint distribution on $(k_i, k_I)$ and not a pair of marginals. In the lagged reading of Section~\ref{sec:cross_scale_shadows} each summand is the divergence of a constituent's belief from its predictive prior. In a functioning participatory system we expect $\mathcal{I}_{\ell \to \ell+1}$ to be small but nonzero for all $\ell < \ell_{\max}$: zero indicates constituents and meta-agent in exact agreement after transport, pointwise consensus (Section~\ref{sec:epistemic_collapse}), while large values indicate a meta-agent that has decoupled from its constituents. The quantity is distinct from the change in the alignment portion of $\mathcal{F}$ when the meta-agent is inserted: that change is signed, whereas $\mathcal{I}_{\ell \to \ell+1}$ is the nonnegative sum of the added child-parent divergences alone.

The complementary top-down channel can be quantified in the same way. Writing the active parent-child relations as a dynamically evolving graph $\mathcal{G} = \{(i, I) : i \in I\}$ whose edges are the channels along which a meta-agent sets a constituent's prior, the per-agent diagnostic
\begin{equation}
\Delta p_i(n) = \int_{\mathcal{C}} D_{\mathrm{KL}}\big(p_i^{(\ell)}(c; n) \big\| p_i^{(\ell)}(c; n-1)\big) \mu_{\mathcal{C}}(dc),
\label{eq:topdown_influence}
\end{equation}
with $\mu_{\mathcal{C}}$ the declared background density on $\mathcal{C}$ and $\Delta p_i > \theta_{\min}$ flagging active top-down feedback, measures how far the evolved parent has moved the constituent's predictive prior between successive episodes.

The construction also makes structural predictions about meta-agents at level $\ell+1$ relative to level $\ell$. In the mean-sector uniform-weight case the meta-agent law is a fixed Markov image of its constituents' laws, and wherever the meta-agent's frame introduces no horizontal defect, the meta-agent's pulled-back information geometry is dominated by its constituents', with equality exactly along directions whose fine score the meta-agent retains~\cite[Ch.~B2, the coarse or meta-agent perceived-geometry corollary]{DennisGaugeVFESuppB2026}. A horizontal defect can reverse the inequality~\cite[Ch.~B2, the exact signed positivity criterion]{DennisGaugeVFESuppB2026}, so the prediction that meta-agent beliefs vary more smoothly over $\mathcal{C}$ than constituent beliefs, with short-wavelength structure averaged out at higher levels, is conditional on frame compatibility, whose exact form is an isotropy condition; this comparison involves base directions and belongs to the base-transport work deferred to Supplement~B. Meta-agent covariances follow the convex-combination pooling of Eq.~\eqref{eq:meta_agent_sigma_impl}, so the construction predicts a parent covariance comparable to the weighted average of the transported constituent covariances, with the strict barycenter adding a positive-semidefinite dispersion term that can only enlarge it. A more confident parent would follow only from the precision-weighted product-of-experts pooling that the meta-agent state construction declines to adopt (Section~\ref{sec:meta_agent_state_construction}). Each prediction requires a multi-seed measurement before any emergence claim can be made (Section~\ref{sec:participatory}).

\section{Coarse-Graining IV: The Ouroboros Tower and Its Declared Root}\label{sec:ouroboros}

Assembling the rungs of the preceding sections yields the Ouroboros tower, the multi-level structure in which each level is a full free-energy functional and adjacent levels are coupled by the cross-scale shadows of Eq.~\eqref{eq:cross_scale_shadow}, read as lagged predictive priors. Wheeler's participatory universe envisions observer-participants whose registrations give tangible reality to what is observed~\cite{Wheeler1983, Wheeler1990}. The tower gives a formal counterpart of one strand of that picture (W4): when coherent clusters are licensed (Section~\ref{sec:meta_agent_variational}; status C), they form meta-agents that set the priors of their constituents. Agents are organized by level $\ell \in \mathbb{N}$: level-0 base agents cohere into level-1 meta-agents, which cluster into level-2 meta-meta-agents, and so on. Each level-$\ell$ agent is a full section $\sigma_I^{(\ell)}$ carrying a primitive belief $q_I^{(\ell)}$, a primitive generative model $s_I^{(\ell)}$, and a frame $\phi_I^{(\ell)}$, with its prior and hyper-prior given by the cross-scale shadows of the level-$(\ell+1)$ meta-agent through Eq.~\eqref{eq:cross_scale_shadow}. The tower is capped at a finite $\ell_{\max}$, and on a finite network a cap is forced in any case (Section~\ref{sec:rg_status}). The reading on which the resulting meta-agents are agents on a par with their constituents, not summaries of them, is an interpretation (status S); the derivation of Section~\ref{sec:meta_agent_variational} supplies instead the reading of the parent as a fitted shared prior.

Constituents receive their shadow from the immediate parent alone. A normalized hierarchy is depth-Markov: each level's factor reads the complete next-level state and no deeper variable, and for an adaptive tower, generated bottom-up by the coarse-graining maps, this holds automatically by Bayes reversal~\cite[Ch.~12]{DennisGaugeVFE2026}. A direct grandparent term is then not a factor of the tower law, and multiplying it in counts information already transmitted through the parent a second time; in a separately declared top-down hierarchy it would require its own normalized kernels. We therefore keep only the depth-one model coupling,
\begin{equation}
\mathcal{F}_{\mathrm{ouro}}\big[s_i; s_I^{(\ell+1)}\big] = \lambda_h D_{\mathrm{KL}}\big(s_i \big\| \tilde\Omega_{iI}\big[s_I^{(\ell+1)}\big]\big),
\label{eq:ouroboros_F}
\end{equation}
which is the model hyper-prior term of Eq.~\eqref{eq:free_energy_functional_final} with $r_i = \tilde\Omega_{iI}[s_I^{(\ell+1)}]$, the depth-one model shadow; deeper ancestors act on $s_i$ only through the chain of depth-one couplings. The cross-scale belief coupling is carried by the self-term $D_{\mathrm{KL}}(q_i \| p_i)$ of Eq.~\eqref{eq:free_energy_functional_final}, whose prior $p_i$ is the depth-one belief shadow $\Omega_{iI}[q_I^{(\ell+1)}]$, and is not re-added here.

A genuine tower predicts that meta-agents at level $\ell+1$ evolve more slowly than their constituents at level $\ell$. This cross-scale slowdown, $t^{(\ell)} < t^{(\ell+1)}$ at matched field type, is a conjecture (status C), not an imposed split. It rests on two mechanisms: meta-agents integrate over constituent beliefs, so high-frequency components of constituent drift are averaged out at the parent level, and changing a meta-agent state requires coherent change across several constituents through the gauge-covariant barycenter, which introduces consensus energy barriers that resist rapid updates. Parent stiffness supplies no third mechanism: with row-normalized $\beta$ the incoming relational stiffness is a convex combination of constituent precisions and does not grow with $|I|$ under the adopted covariance pooling. The within-level separation $t_{\mathrm{belief}}^{(\ell)} < t_{\mathrm{model}}^{(\ell)}$, by contrast, is imposed by the adiabatic hypothesis of Section~\ref{sec:timescales}, in which beliefs and attention rows form the fast block and models and prior parameters the slow block~\cite[Ch.~8]{DennisGaugeVFE2026}. The fully interleaved ordering $t_{\mathrm{model}}^{(\ell)} < t_{\mathrm{belief}}^{(\ell+1)}$ that a single chain would assert is supported by no mechanism here and is in tension with the prior-inheritance structure, under which level-$\ell$ priors inherit their timescales from the parent. The cross-scale slowdown is analogous to collective variables equilibrating more slowly than microscopic ones and to institutional change being slower than individual opinion change. It is not built into the construction, which uses one step size at every level; testing it requires a declared per-level timescale measurement.

The tower terminates at its top level. There is no parent meta-agent, and the top-level hyper-prior $r_i^{(\ell_{\max})}$ is a declared prior, the boundary condition of the hierarchy. For the belief channel an energy-level closure is available: top-level agent $i$ may take as its prior the coherence-weighted pool of the transported beliefs of all active agents at all levels. Mixtures of transported Gaussians leave the Gaussian family on which the closed-form divergences rely, so the pool is taken as the Gaussian moment projection of the mixture $\bar Q_i = \sum_j \bar w_j (\Omega_{ij})_{\#} q_j$, the minimizer of $D_{\mathrm{KL}}(\bar Q_i \| p)$ over Gaussians $p$,
\begin{equation}
\mu_i^{(\mathrm{top})} = \sum_j \bar w_j \Omega_{ij}\mu_j,
\qquad
\Sigma_i^{(\mathrm{top})} = \sum_j \bar w_j \Big[\Omega_{ij} \Sigma_j \Omega_{ij}^\top + \big(\Omega_{ij}\mu_j - \mu_i^{(\mathrm{top})}\big)\big(\Omega_{ij}\mu_j - \mu_i^{(\mathrm{top})}\big)^\top\Big],
\end{equation}
with normalized weights $\bar w_j = w_j / \sum_k w_k$ and $w_j \propto \exp(-\bar D_j)$, where $\bar D_j$ is agent $j$'s mean transported belief divergence from the collective, so that the weights favor agents coherent with the collective. The between-mean dispersion term is retained: dropping it, as the working form Eq.~\eqref{eq:meta_agent_sigma_impl} of the cluster barycenter does, would depart from the moment projection, and because the apex pool is not gated by coherence, the dispersion term is what records a globally incoherent collective as a broad prior. The model channel is not pooled.

The apex closes the tower at the level of the consensus energy and its update schedule, not as a generative model. Every finite acyclic generative model in which each record has a latent parent leaves at least one latent whose prior is declared and not constituted by the agents it organizes~\cite[Ch.~12, the declared-root proposition]{DennisGaugeVFE2026} (status P), and a reciprocal same-time structure is a probability model only as a globally normalized joint with its partition function~\cite[Ch.~4]{DennisGaugeVFE2026}. The pooled prior reads the current beliefs of the agents it organizes, so it is not a factor of one joint. The declared root is the formal residue of Wheeler's ``no tower of turtles''~\cite{Wheeler1990}, a reading of status S: the regress of priors stops at a declared prior and does not close on itself, and the companion theory accordingly replaces a self-closing apex by the declared root. The participatory content survives in three typed forms, each a sequence of fixed normalized joints~\cite[Ch.~12]{DennisGaugeVFE2026}: the inference loop, in which parent and child updates read each other under block coordinate ascent, as in any connected hierarchical model; slow learning of a declared top prior, whose optimum over an unrestricted family is the aggregate of the recognition laws it summarizes and, within the Gaussian family, their moment match including the dispersion retained above; and delayed feedback, in which the parent inferred at one episode sets its children's prior at the next (Section~\ref{sec:cross_scale_shadows}). The apex pool is the same-time, energy-level counterpart of the second: the companion theory aggregates the recognition laws of one top latent across episodes by a slow parameter update, whereas the apex pools across agents of all levels at one time. It has the feedback structure of Wheeler's self-excited circuit without being derived from it or evidence for it (status S). It is not that circuit, which runs through time, is constituted by registrations, and sustains itself~\cite{Wheeler1983}: the dynamics here relax to stationarity unless observations are injected, and the top-level model channel is a fixed boundary condition.

A participatory system must keep evolving. Stationarity of the specified dynamics means stationarity of $\mathcal{F}$, at which beliefs, priors, and frames cease evolving. Generic stationary points balance alignment against prior attachment, observation fit, and the attention entropy, and therefore retain persistent disagreement, so perfect consensus is one special stationary configuration and not the generic one; what makes a stationary state non-participatory is the cessation of evolution itself. The notion is non-stationarity of a deterministic descent. It is not the statistical-physics sense of distance from equilibrium, which refers to broken detailed balance or entropy production in a driven stochastic dynamics; the dynamics here is relaxational. Three indicators track it, each rendered dimensionless against a declared early-run reference window before any combination: an energy flux $\Phi_E(t) \propto |d\mathcal{F}/dt|$, an information flux $\Phi_I(t)$ tracking the magnitude of entropy change across belief and prior, and a gradient variance $V_{\nabla}(t) \propto \mathrm{Var}_{i}(\|\nabla_{q_i}\mathcal{F}\|^2)$ taken across agents at fixed $t$. They combine into the non-stationarity score $S_{\mathrm{ns}}(t) = (\Phi_E + \Phi_I + V_{\nabla})/3$, whose value depends on the reference window; values such as $0.5$ and $1$ are candidate thresholds to preregister and calibrate, not validated ones. Without observation drive the flow relaxes. The reduced functional is bounded below and real-analytic in the Gaussian parameters, so by the {\L}ojasiewicz gradient inequality~\cite{Lojasiewicz1963} every natural-gradient trajectory with precompact image converges to a stationary point and has finite total Fisher length~\cite[Ch.~7, the finite-duration proposition]{DennisGaugeVFE2026}; the corresponding statement for discrete descent iterates is~\cite{absil2005convergence}. Of the mechanisms that might sustain evolution, only observation injection supplies persistent exogenous drive. Heterogeneous top-down feedback and spatial separation shape how the system relaxes without preventing relaxation, and frame differences cannot maintain disagreement by themselves, since Regime~I frames are removable at a point. Whether observation injection sustains non-stationarity is open (status C).

Every per-level free-energy term is integrated over $\mathcal{C}$ against the same declared background density, the volume form of the background chart (Section~\ref{sec:flat_base}). That density is an input, not the geometry the agents induce. The induced geometry, the pulled-back metrics of Section~\ref{sec:pullback}, enters no free-energy term, so the participation loop of the tower closes only within information, from belief to meta-belief to predictive prior.

\section{Prospective Validation of Multi-Scale Formation}\label{sec:participatory}

The construction of Sections~\ref{sec:meta_agent_formation}--\ref{sec:ouroboros} defines how coherent clusters may form meta-agents and how their beliefs and models may propagate across levels. No simulation result is claimed (Section~\ref{sec:scope_limitations}); this section states the experiment that would test the multi-level claims. The formation rule under test is the threshold detector and barycenter of Section~\ref{sec:meta_agent_formation}; the companion theory replaces both by normalized marked parents~\cite[Ch.~13]{DennisGaugeVFE2026}, which is the rule a later protocol should test.

Formation alone cannot discriminate participation from seeding. On the homogeneous four-agent cycle of the companion theory, a collective forms from every Gaussian initial law with probability bounded below uniformly in the regime parameter, whether or not the source prior learns, because the random initial source assignment seeds the pair~\cite[Ch.~14, the seeded-contrast-formation proposition]{DennisGaugeVFE2026}. The source prior there is the attention prior $\pi_{ij}$ of this paper. Stable multi-child clusters, reproduction across seeds, and an active model channel would therefore all be met by a system whose attention prior never learns. The decidable form of the participatory reading is therefore the companion theory's pairing-selection conjecture (status C): slow learning of the attention prior selects which pair forms a collective, at the level of the law, against a frozen-prior control whose pairing law stays symmetric~\cite[Ch.~14, the pairing-selection conjecture]{DennisGaugeVFE2026}. The protocol below is its finite-sample counterpart in the present construction; neither bears on the participatory-physics readings of Part~II.

The protocol has three arms at one declared configuration, differing only as named. Arm A has a learning attention prior and an active, untied model channel with a declared nonzero $\gamma_{ij}$, so that the term $\gamma_{ij} D_{\mathrm{KL}}(s_i\|\tilde\Omega_{ij}[s_j])$ is exercised and not inferred from a frozen or tied channel. Arm B, the frozen-attention-prior control, is identical except that $\pi_{ij}$ does not learn. Arm C is identical to arm A except that the model channel is frozen ($\lambda_h = 0$ and $\gamma_{ij} = 0$). The arms share seeds: each seed fixes the initial means and the initial assignment draws for all three arms, so that comparisons between arms are paired.

The following quantities are declared here and may be revised only before execution, with any revision reported. The formation event at a recorded step is the existence of a connected block $I$ with $2 \le |I| \le N-1$ whose internal belief and model transported relative entropies are all at most $\delta_{\mathrm{in}} = 10^{-2}$ nats and whose boundary belief transported relative entropies all exceed $\delta_{\mathrm{out}} = 1/2$ nat, the form of the event in the companion proposition with its internal bound relaxed. The persistence window has flow-time length $W = 10/\eta$, ten relaxation times of an isolated belief mode under the natural-gradient flow of rate $\eta$, and ends at the declared horizon. The primary statistic for each seed and arm is the indicator $H$ that one block satisfies the formation event throughout the window, together with the identity of that block. Each arm runs $n = 20$ seeds. The formation claim passes in arm A if $H = 1$ in at least $16$ of the $20$ seeds. The participatory claim requires a difference between arms A and B: it passes only if an exact two-sided McNemar test on the paired indicators rejects at level $0.05$, or if an exact permutation test at level $0.05$ on the paired block identities rejects equality of the laws of the formed blocks. With twenty paired seeds no single aberrant seed decides the outcome, and the McNemar test needs at least six discordant pairs in one direction ($2 \cdot 2^{-6} = 0.031$). The model-channel claim passes if arms A and C differ by the same McNemar criterion.

The formation claim is refuted for the tested configuration if arm A fails its pass rule. The participatory claim is unsupported if arms A and B do not differ by the pass rule, whatever arm A shows. A result that appears only with the model channel frozen, only under thresholds chosen after inspection, or only in one seed validates nothing. Two conditions are preconditions for a valid measurement and not refutation criteria: population accounting must close at every recorded step, $N_{\mathrm{total}}(t) = \sum_\ell N_\ell(t)$ with $N_\ell(t)$ the number of active agents at level $\ell$, with every meta-agent's constituents and parent recorded; and the raw per-step states (beliefs, models, priors, hyper-priors, frames, transported pairwise divergences, $\beta_{ij}$, $\gamma_{ij}$, detector scores, assignments, energy components, and cross-scale links) must be retained so that every reported statistic can be recomputed. A failure of either invalidates the run as a measurement; it does not refute the tower.

Until this protocol is completed, the Ouroboros tower supplies a mathematical construction and a set of falsifiable predictions (status C), not evidence that the dynamics generate a persistent hierarchy or sustained non-stationarity.

